\chapter{Introduction}
\label{ch:intro}

\section{Background and motivation}

Automatic topic classification of news text is one of the oldest and most durable applications of
natural language processing. News agencies route incoming wire copy to desks; aggregators assemble
sections; recommendation systems filter by interest; archives index by subject. In each case the
underlying computational problem is the same: given a short passage of text, assign it to one of a
small set of editorial categories.

What has changed repeatedly over the past three decades is not the problem but the representation.
Early systems treated a document as an unordered bag of words and applied a linear classifier to
sparse count features. The introduction of distributed word representations replaced sparse
indicators with dense vectors that encoded distributional similarity. Recurrent neural networks
then made it possible to model a document as a \emph{sequence} rather than a set, processing tokens
in order and maintaining a hidden state that in principle summarises everything seen so far. The
Transformer subsequently discarded recurrence altogether, replacing the sequential state with an
attention mechanism that relates every position to every other position directly.

The Transformer's dominance in contemporary NLP is not in dispute. What is less often examined
carefully is \emph{where} its advantage actually comes from, and whether that advantage is present
in every task to which the architecture is applied. The mechanism that makes self-attention
powerful --- a constant-length computational path between any two positions in a sequence --- is
explicitly a statement about \emph{long-range} dependency. It is reasonable to ask what remains of
that advantage when the sequences in question are not long.

This question is not merely academic. A great many deployed text classification systems operate on
short inputs: search queries, product titles, support ticket subjects, headlines. If the
architectural advantage of self-attention is contingent on sequence length, then the
engineering decision for such systems cannot be read off from benchmark leaderboards built on
longer documents.

\section{Problem statement}

The task addressed in this report is four-class topic classification of short news documents. Given
a news item consisting of a headline and a one-sentence description, the system must assign exactly
one of four labels: \emph{World}, \emph{Sports}, \emph{Business}, or \emph{Science/Technology}.

Two properties of this task make it a suitable vehicle for the architectural comparison at hand.

First, the documents are short. In the corpus used here the median document is 39 tokens long. This
places the task squarely in the regime where the long-range-dependency argument for self-attention
is weakest, and therefore where a direct comparison is most informative.

Second, the label set has four classes rather than two. A binary sentiment task yields a single
accuracy number and a largely uninformative confusion structure. Four classes produce a confusion
matrix with twelve off-diagonal cells, which permits a far more detailed account of \emph{how} a
model fails rather than merely how often. As will be shown in Chapter~\ref{ch:results}, the error
structure turns out to be considerably more informative than the headline accuracy.

\section{Research question}

The report addresses a single primary research question:

\begin{quote}
\textbf{How do recurrent networks and Transformers compare at classifying short news texts?}
\end{quote}

This decomposes into four subsidiary questions, each of which is answered empirically:

\begin{enumerate}
\item \textbf{Accuracy.} Does either architecture achieve higher classification accuracy, and is
      any observed difference statistically distinguishable from sampling noise?
\item \textbf{Error structure.} Do the two architectures fail on the same documents and the same
      class boundaries, or do they exhibit different inductive biases?
\item \textbf{Cost.} How do the architectures compare in training time, inference latency and
      parameter efficiency at matched capacity?
\item \textbf{Explanation.} Whatever differences or equivalences are observed, what property of the
      task accounts for them?
\end{enumerate}

\section{Objectives}

The study was conducted with the following concrete objectives.

\begin{enumerate}
\item To implement a bidirectional LSTM classifier from first principles, writing the gating
      equations, the cell-state recurrence and the bidirectional unrolling explicitly rather than
      invoking a library primitive.
\item To implement a Transformer encoder classifier from first principles, writing scaled
      dot-product attention, multi-head projection and merging, sinusoidal positional encoding,
      layer normalisation and the residual block structure explicitly.
\item To verify both implementations numerically against reference implementations and against
      invariants that any correct implementation must satisfy, \emph{before} drawing any
      experimental conclusion from them.
\item To design a training and evaluation protocol in which architecture is the only free
      variable, holding data, preprocessing, capacity and training budget fixed.
\item To evaluate both models on a held-out test set using metrics appropriate to a balanced
      multi-class problem, including a test of statistical significance on the paired predictions.
\item To analyse the computational cost of both architectures and to interpret the result in terms
      of the structural properties of recurrence and attention.
\end{enumerate}

\section{Scope and delimitations}

It is as important to state what this study does \emph{not} attempt as what it does.

\keypoint{Training from scratch, not transfer learning.} Both models are trained from randomly
initialised embeddings on the task data alone. No pretrained representations are used. This is a
deliberate scoping decision: it isolates the architectural question. A comparison involving a
pretrained Transformer would primarily measure the value of pretraining, which is a different and
already well-established result.

\keypoint{Matched capacity, not maximum performance.} The models are deliberately sized to have
near-identical non-embedding parameter counts. Neither is tuned to its maximum achievable accuracy
on AG News. The objective is a fair comparison, not a leaderboard entry.

\keypoint{A controlled subset, not the full corpus.} Following the case-study brief, 60{,}000 of
the 120{,}000 available training records are used. The consequences of this choice are discussed
in Section~\ref{sec:threats}.

\keypoint{One corpus, one language, one document length regime.} All conclusions are drawn from
English-language short news text. The report is explicit, particularly in
Chapter~\ref{ch:discussion}, about which findings are expected to generalise and which are
specific to this regime.

\section{Contributions}

This report makes the following contributions.

\begin{enumerate}
\item \textbf{A fully transparent implementation} of both architectures, in which every
      sequence-modelling operation is written explicitly and is available for inspection. Nothing
      of consequence is hidden inside a library call.

\item \textbf{A verification methodology for hand-written layers.} Nine numerical checks are
      specified, motivated and executed, covering equivalence against reference implementations,
      mathematical invariants, masking correctness and gradient viability. The report argues that
      such verification is not optional for from-scratch work, because the characteristic failure
      mode is silent degradation rather than an observable error.

\item \textbf{A genuinely controlled comparison}, in which data, preprocessing, capacity and
      training budget are held fixed and only the architecture varies.

\item \textbf{A statistically disciplined negative result.} The observed accuracy difference is
      0.09 percentage points. Rather than reporting a winner, the study applies McNemar's exact
      test, finds $p = 0.754$, and concludes that the architectures are indistinguishable on this
      task. The report treats the significance test as essential rather than decorative.

\item \textbf{A cost finding that contradicts the expected result}, together with an analysis that
      distinguishes the asymptotic argument for attention from the constant factors that determine
      wall-clock time at short sequence lengths.
\end{enumerate}

\section{Organisation of this report}

The remainder of the report is structured as follows.

\textbf{Chapter~\ref{ch:literature}} reviews the development of text classification from sparse
feature methods through recurrent architectures to self-attention, and surveys published results
on the AG News benchmark in order to situate the present work.

\textbf{Chapter~\ref{ch:theory}} develops the mathematical background: the recurrent formulation
and its gradient pathology, the LSTM gating solution, the attention mechanism and its scaling, and
a complexity analysis that formalises the structural trade-off between the two architectures.

\textbf{Chapter~\ref{ch:method}} describes the corpus, the data cleaning required, the controlled
experimental design, the preprocessing pipeline and the evaluation metrics.

\textbf{Chapter~\ref{ch:implementation}} presents the implementation of both models, with
attention to the specific points at which a from-scratch implementation is most likely to go
wrong.

\textbf{Chapter~\ref{ch:verification}} presents the nine verification checks, their motivation and
their results.

\textbf{Chapter~\ref{ch:results}} reports the experimental results: training dynamics, accuracy,
per-class performance, confusion structure, statistical significance, computational cost and a
qualitative error analysis.

\textbf{Chapter~\ref{ch:discussion}} interprets the results, answers the research question,
compares the findings with published work, and states the threats to validity.

\textbf{Chapter~\ref{ch:conclusion}} concludes and identifies future work.

The appendices contain the full source listings of the hand-written components, the complete
classification reports, and instructions for reproducing every number in this document.
